\section{Теоретические основы}
\label{sec:Theory} \index{Theory}

В данной главе вводятся теоретические сведения, необходимые для понимания предлагаемого решения: архитектура и режим генерации больших языковых моделей, парадигма пошаговых рассуждений, методы обучения с подкреплением применительно к LLM, формализм графов знаний и мера структурного расстояния между графами, а также параметр-эффективное дообучение.

\subsection{Большие языковые модели и авторегрессионная генерация}
\label{subsec:theory-llm}

Современные большие языковые модели~\cite{brown2020language, touvron2023llama} строятся на архитектуре \textit{трансформера}, ключевым элементом которой является механизм многоголового самовнимания (multi-head self-attention). Для последовательности скрытых представлений $X \in \mathbb{R}^{n \times d}$ операция внимания вычисляется как
\begin{align}
    \mathrm{Attention}(Q, K, V) = \mathrm{softmax}\!\left(\frac{QK^\top}{\sqrt{d_k}}\right) V,
    \label{eq:attention}
\end{align}
где $Q = XW_Q$, $K = XW_K$, $V = XW_V$~--- проекции запросов, ключей и значений, а $d_k$~--- размерность ключа. Самовнимание позволяет каждому токену агрегировать информацию из всего контекста, что лежит в основе способности модели улавливать долгосрочные зависимости.

Языковая модель задаёт распределение над последовательностями токенов $y = (y_1, \ldots, y_T)$ через авторегрессионную факторизацию
\begin{align}
    \pi_\theta(y \mid x) = \prod_{t=1}^{T} \pi_\theta(y_t \mid x, y_{<t}),
    \label{eq:autoregressive}
\end{align}
где $x$~--- входной запрос (prompt), $y_{<t}$~--- ранее сгенерированные токены, $\theta$~--- параметры модели. Генерация осуществляется последовательно: на каждом шаге из распределения $\pi_\theta(\cdot \mid x, y_{<t})$ выбирается очередной токен (жадно, сэмплированием с температурой $T$ или другими стратегиями декодирования). Именно представление политики в виде распределения~(\ref{eq:autoregressive}) делает возможным применение методов обучения с подкреплением, рассматривающих генерацию текста как последовательность действий.

\subsection{Цепочки рассуждений}
\label{subsec:theory-cot}

Метод Chain-of-Thought (CoT)~\cite{wei2022chain} показал, что побуждение модели генерировать промежуточные шаги рассуждения перед выдачей финального ответа существенно повышает точность на задачах, требующих логического и арифметического вывода. Формально вместо прямого моделирования $\pi_\theta(a \mid q)$ модель порождает промежуточную цепочку $z$ (rationale) и ответ $a$:
\begin{align}
    \pi_\theta(a, z \mid q) = \pi_\theta(z \mid q)\, \pi_\theta(a \mid q, z).
\end{align}
Развитием CoT являются Self-Consistency~\cite{wang2023selfconsistency} (агрегация ответов по множеству сэмплированных цепочек), Tree-of-Thought~\cite{yao2023tree} (поиск по дереву частичных рассуждений) и Graph-of-Thought~\cite{besta2024graph} (представление рассуждения произвольным графом). Общим ограничением этих методов является то, что они действуют исключительно на этапе инференса и не оценивают \textit{структурную} корректность промежуточных шагов~--- именно этот пробел адресует настоящая работа.

\subsection{Обучение с подкреплением для языковых моделей}
\label{subsec:theory-rl}

В парадигме обучения с подкреплением генерация ответа $y$ на запрос $x$ рассматривается как эпизод, в котором политика $\pi_\theta$ получает скалярную награду $R(x, y)$. Цель обучения~--- максимизация ожидаемой награды с регуляризацией отклонения от опорной политики $\pi_{\mathrm{ref}}$:
\begin{align}
    \max_{\theta}\ \mathbb{E}_{x \sim \mathcal{D},\, y \sim \pi_\theta(\cdot \mid x)} \big[ R(x, y) \big] - \beta\, \mathrm{KL}\!\left[\pi_\theta(\cdot \mid x)\,\|\,\pi_{\mathrm{ref}}(\cdot \mid x)\right].
    \label{eq:rl-objective}
\end{align}
KL-член предотвращает чрезмерное удаление политики от опорной модели и тем самым стабилизирует обучение.

\textbf{Policy gradient и REINFORCE.} Базовым результатом, лежащим в основе методов оптимизации политики, является теорема о градиенте политики, дающая несмещённую оценку градиента ожидаемой награды:
\begin{align}
    \nabla_\theta\, \mathbb{E}_{y \sim \pi_\theta}\big[R(x, y)\big] = \mathbb{E}_{y \sim \pi_\theta}\big[R(x, y)\, \nabla_\theta \log \pi_\theta(y \mid x)\big].
    \label{eq:policy-gradient}
\end{align}
Алгоритм REINFORCE напрямую использует~(\ref{eq:policy-gradient}), однако обладает высокой дисперсией оценки. Для её снижения вводят базовую линию (baseline) $b(x)$, вычитаемую из награды: $\hat{A} = R(x, y) - b(x)$. Современные методы (PPO, GRPO) развивают эту идею, используя более совершенные оценки преимущества и ограничивая величину шага обновления политики.

\textbf{PPO.} Алгоритм Proximal Policy Optimization, лежащий в основе классического RLHF~\cite{ouyang2022training}, оптимизирует клиппированную суррогатную функцию
\begin{align}
    \mathcal{L}^{\mathrm{PPO}}(\theta) = \mathbb{E}_t\!\left[\min\!\left(\rho_t(\theta)\, \hat{A}_t,\ \mathrm{clip}(\rho_t(\theta),\, 1-\varepsilon,\, 1+\varepsilon)\, \hat{A}_t\right)\right],
    \label{eq:ppo}
\end{align}
где $\rho_t(\theta) = \frac{\pi_\theta(y_t \mid x, y_{<t})}{\pi_{\theta_{\mathrm{old}}}(y_t \mid x, y_{<t})}$~--- отношение вероятностей, $\hat{A}_t$~--- оценка преимущества, $\varepsilon$~--- параметр клиппинга. Оценка $\hat{A}_t$ обычно требует обучения отдельной value-сети (критика), что удваивает потребление памяти.

\textbf{GRPO.} Алгоритм Group Relative Policy Optimization~\cite{shao2024deepseekmath} устраняет необходимость в критике: для каждого запроса $x$ сэмплируется группа из $G$ ответов $\{y^{(1)}, \ldots, y^{(G)}\}$, и преимущество каждого ответа оценивается через нормировку наград внутри группы:
\begin{align}
    \hat{A}^{(i)} = \frac{R(x, y^{(i)}) - \mathrm{mean}\{R(x, y^{(j)})\}_{j=1}^{G}}{\mathrm{std}\{R(x, y^{(j)})\}_{j=1}^{G}}.
    \label{eq:grpo-advantage}
\end{align}
Это снижает требования к памяти примерно вдвое и делает алгоритм естественно пригодным для разреженных и композитных наград, что критично для задачи настоящей работы (раздел~\ref{subsec:reward-design}).

\subsection{Графы знаний и расстояние редактирования графов}
\label{subsec:theory-kg}

\textit{Граф знаний}~\cite{vrandecic2014wikidata} представляет факты в виде троек $(s, r, o)$, где $s$~--- субъект, $r$~--- отношение, $o$~--- объект. Совокупность троек образует ориентированный мультиграф, вершинами которого выступают сущности, а рёбрами~--- отношения. Крупномасштабные графы (Wikidata) содержат миллионы сущностей, образуя фактическую основу для проверки корректности рассуждений.

Для количественного сравнения двух графов используется \textit{расстояние редактирования графов} (Graph Edit Distance, GED)~--- минимальная суммарная стоимость элементарных операций (вставка, удаление, замена вершин и рёбер), переводящих один граф в другой:
\begin{align}
    \mathrm{GED}(G_1, G_2) = \min_{(e_1, \ldots, e_k) \in \mathcal{P}(G_1, G_2)} \sum_{i=1}^{k} c(e_i),
    \label{eq:ged}
\end{align}
где $\mathcal{P}(G_1, G_2)$~--- множество последовательностей операций, переводящих $G_1$ в $G_2$, а $c(\cdot)$~--- стоимость операции. Вычисление GED в общем случае NP-трудно, однако для малых графов (единицы~--- десятки вершин) выполнимо точными алгоритмами. На основе GED строится нормированная мера сходства, используемая в работе (формула~(\ref{eq:ges})).

\subsection{Метрики точности, полноты и F1-меры}
\label{subsec:theory-metrics}

Для оценки соответствия двух множеств (например, сгенерированных и эталонных вершин или рёбер графа) используются классические метрики информационного поиска. Пусть $A$~--- множество предсказанных элементов, $B$~--- множество эталонных. Тогда
\begin{align}
    \mathrm{Precision} = \frac{|A \cap B|}{|A|}, \qquad
    \mathrm{Recall} = \frac{|A \cap B|}{|B|}, \qquad
    F_1 = \frac{2\,\mathrm{Precision}\cdot\mathrm{Recall}}{\mathrm{Precision} + \mathrm{Recall}}.
    \label{eq:prf}
\end{align}
Точность (precision) характеризует долю корректных элементов среди предсказанных, полнота (recall)~--- долю покрытых эталонных элементов, а $F_1$-мера~--- их гармоническое среднее, штрафующее дисбаланс между точностью и полнотой. Эти метрики применяются в работе на уровне вершин и рёбер графа рассуждения (определение~\ref{def:metrics}); их сочетание позволяет различать модели, которые генерируют много лишних троек (высокий recall, низкая precision), и модели, выдающие лишь часть корректной структуры (высокая precision, низкий recall).

\subsection{Параметр-эффективное дообучение}
\label{subsec:theory-lora}

Полное дообучение всех параметров крупной модели требует значительных вычислительных ресурсов. Метод Low-Rank Adaptation (LoRA)~\cite{hu2022lora} замораживает исходные веса $W_0 \in \mathbb{R}^{d \times k}$ и обучает лишь низкоранговую поправку:
\begin{align}
    W = W_0 + \Delta W = W_0 + BA, \qquad B \in \mathbb{R}^{d \times r},\ A \in \mathbb{R}^{r \times k},\ r \ll \min(d, k).
    \label{eq:lora}
\end{align}
Поскольку ранг $r$ мал, число обучаемых параметров сокращается до доли процента от исходного, что позволяет дообучать модель на одной GPU. LoRA-адаптеры обычно применяются к проекционным матрицам внимания ($W_Q, W_K, W_V$) и слоям MLP. В сочетании с пейджингом состояния оптимизатора и вычислениями в формате \texttt{bfloat16} это делает RL-дообучение компактных моделей практически доступным.
